\subsection{与正则局部环关联的分次环}

定理 1. --- 设 A 为诺特局部环。下列条件彼此等价：

\begin{enumerate}
\def\labelenumi{(\roman{enumi})}
\item
  A 是正则的。
\item
  理想 \(\mathfrak{m}_{\mathbf{A}}\) 由 A 的一个正割子集生成（§ 3, no. 2, 定义 1）。
\item
  理想 \(\mathfrak{m}_{\mathrm{A}}\) 由 A 的一个完全正割序列生成（A, X, p.~157, 定义 2）。
\item
  令 S 为 \(\kappa_{\mathbf{A}}\)-向量空间 \(\mathfrak{m}_{\mathbf{A}} / \mathfrak{m}_{\mathbf{A}}^{2}\) 的对称代数，并令 \(\operatorname{gr}(\mathbf{A}) = \bigoplus_{n \geqslant 0} \mathfrak{m}_{\mathbf{A}}^{n} / \mathfrak{m}_{\mathbf{A}}^{n+1}\) 为与 A 关联的分次环。S 的典范同态 \(\gamma\) 映到 \(\operatorname{gr}(\mathbf{A})\) 上是双射。
\item
  存在整数 \(r \geqslant 0\)，使得 \(H_{\mathbf{A},\mathfrak{m}_{\mathbf{A}}} = (1 - \mathrm{T})^{-r}\)，即 \(\mathfrak{m}_{\mathbf{A}} = 0\) 当 \(r = 0\) 时；并且等式 \([\mathfrak{m}_{\mathbf{A}}^{n}/\mathfrak{m}_{\mathbf{A}}^{n+1}:\kappa_{\mathbf{A}}] = \binom{n + r - 1}{r - 1}\) 对每个满足 \(n \geqslant 0\) 的整数成立（在 \(r > 0\) 时）。
\item
  我们有 \(\mathrm{H}_{\mathrm{A},\mathfrak{m}_{\mathrm{A}}} = (1 - \mathrm{T})^{-d}\)，其中 \(d = \dim (\mathrm{A})\)。
\end{enumerate}

如果这些条件成立，则 A 的每个坐标系都是 A 的一个完全正割序列。

\begin{enumerate}
\def\labelenumi{(\roman{enumi})}
\setcounter{enumi}{1}
\item
  \(\Rightarrow\) (i)：事实上，每个正割子集至多含有 dim(A) 个元素（ \(\S 3\) , no. 2, 定理 1）。
\item
  \(\Rightarrow\) (ii)：事实上，每个完全正割序列都是正割的（§ 3, no. 2, 命题 3 的推论）。
\end{enumerate}

\((\mathrm{iv}) \Rightarrow (\mathrm{iii}) : \text{设 } (x_1, ..., x_r)\) 是 \(\mathfrak{m}_{\mathbf{A}}\) 中元素的一个序列，其模 \(\mathfrak{m}_{\mathbf{A}}^2\) 的类形成 \((\xi_1, ..., \xi_r)\)，后者是 \(\mathfrak{m}_{\mathbf{A}} / \mathfrak{m}_{\mathbf{A}}^2\) 关于域 \(\kappa_{\mathbf{A}}\) 的一个基。若性质 (iv) 成立，则 \(\operatorname{gr}(\mathbf{A})\) 是多项式代数 \(\kappa_{\mathbf{A}}[\xi_1, ..., \xi_r]\)，且序列 \((x_1, ..., x_r)\) 是完全正割序列（A, X, p.~160, 定理 1）。这也证明了定理 1 的最后一个断言。

\begin{enumerate}
\def\labelenumi{(\roman{enumi})}
\tightlist
\item
  \(\Rightarrow\) (iv)：置 \(r = [\mathfrak{m}_{\mathbf{A}} / \mathfrak{m}_{\mathbf{A}}^{2}:\kappa_{\mathbf{A}}]\)。根据 § 4, no. 2 的公式 (6)，分次向量空间 S 关于域 \(\kappa_{\mathbf{A}}\) 的庞加莱级数等于
\end{enumerate}

\[
\mathrm{P} _ {\mathrm{S}} = \sum_ {n \geqslant 0} \left[ \mathrm{S} ^ {n} (\mathfrak {m} _ {\mathrm{A}} / \mathfrak {m} _ {\mathrm{A}} ^ {2}): \kappa_ {\mathrm{A}} \right] \mathrm{T} ^ {n} = (1 - \mathrm{T}) ^ {- r}.
\]

设典范同态 \(\gamma: S \to \mathrm{gr}(A)\) 不是双射。由于 \(\gamma\) 是满射，S 中存在齐次元素 \(u\)，它属于 \(S\)，次数为 \(d > 0\)，且被 \(\gamma\) 消去。于是有

\[
\begin{array}{r l} \mathrm {H_ {A,m_ {A}}} = & \mathrm {P_ {S}} - \symrm {P_ {Ker(\gamma)}} \leqslant \mathrm {P_ {S}} - \mathrm {P_ {uS}} = (1 - \mathrm{T} ^ {d}) / (1 - \mathrm{T}) ^ {r} = \\ & = (1 + \mathrm{T} + \dots + \mathrm{T} ^ {d - 1}) / (1 - \mathrm{T}) ^ {r - 1}. \end{array}
\]

根据 § 4, no. 4 的定理 3 和 § 4, no. 1 的引理 2，因而有 dim(A) \textless{} r，且 A 不是正则的。

最后证明条件 (iv) 至 (vi) 的等价性。现在，(iv) 意味着我们有 \(\mathrm{H}_{\mathrm{A},\mathfrak{m}_{\mathrm{A}}}=(1-\mathrm{T})^{-s}\)，其中 \(s=[\mathfrak{m}_{\mathrm{A}}/\mathfrak{m}_{\mathrm{A}}^{2}:\kappa_{\mathrm{A}}]\)。因此，条件 (iv)、(v)、(vi) 意味着我们有 \(\mathrm{H}_{\mathrm{A},\mathfrak{m}_{\mathrm{A}}}=(1-\mathrm{T})^{-m}\)，其中相应地 \(m=[\mathfrak{m}_{\mathrm{A}}/\mathfrak{m}_{\mathrm{A}}^{2}:\kappa_{\mathrm{A}}]\)、\(m\geqslant0\)、\(m=\dim(\mathrm{A})\)。但是，若有 \(\mathrm{H}_{\mathrm{A},\mathfrak{m}_{\mathrm{A}}}=(1-\mathrm{T})^{-m}\)，则根据 §4, no. 4, 定理 3，有 \(\dim(\mathrm{A})=m\)，并且有 \([\mathfrak{m}_{\mathrm{A}}/\mathfrak{m}_{\mathrm{A}}^{2}:\kappa_{\mathrm{A}}]=m\)（因为 \((1-\mathrm{T})^{-m}=1+m\mathrm{T}+\cdots\)）。条件 (iv) 与 (vi) 的等价性立即得到。

推论 1. --- 每个正则诺特局部环都是整闭的，因而特别是整环。

假设 A 正则。则 gr(A) 同构于域上有限多个不定元的多项式代数（定理 1, (iv)）。因此，gr(A) 是诺特整闭环（V, § 1, no. 3, 命题 13 的推论 3），从而 A 是整闭的（V, § 1, no. 4, 命题 15）。

我们将在后面的章节中看到，每个正则诺特局部环都是唯一分解环。

推论 2.--- 设 A、B 是诺特局部环，σ 是从 B 到 A 的局部同态。设 A 正则且 B 完备。σ 为双射的充要条件是：σ 将 κ \(_{B}\) 双射到 κ \(_{A}\)，并将 m \(_{B}\) /m \(_{B}^{2}\) 双射到 m \(_{A}\) /m \(_{A}^{2}\)。

所述条件显然是必要的。

反之，假设 \(\sigma\) 将 \(\mathrm{gr}_{0}(\mathbf{B})\) 同构到 \(\mathrm{gr}_{0}(\mathbf{A})\)，并将 \(\mathrm{gr}_{1}(\mathbf{B})\) 同构到 \(\mathrm{gr}_{1}(\mathbf{A})\)。由于环 \(\mathrm{gr}(\mathbf{B})\) 由 \(\mathrm{gr}_{0}(\mathbf{B}) \cup \mathrm{gr}_{1}(\mathbf{B})\) 生成，而 \(\mathrm{gr}(\mathbf{A})\) 是 \(\mathrm{gr}_{1}(\mathbf{A})\) 作为 \(\mathrm{gr}_{0}(\mathbf{A})\) 上向量空间的对称代数，故同态 \(\mathrm{gr}(\sigma)\) 是双射。因此，\(\sigma\) 是双射（III, § 2, no. 8, 定理 1 的推论 3）。

推论 3. --- 设 \(k\) 为域，A 为局部诺特 \(k\)-代数，且其剩余域等于 \(k\)。A 正则的充要条件是：其完备化 \(\hat{A}\) 同构于一个 \(k\)-代数，该代数由形式级数 \(k[[X_1, ..., X_n]]\) 构成。

这由定理 1 中 (i) 与 (iv) 的等价性以及 III, § 2, no. 9 的命题 11 得出。

